% ════════════════════════════════════════════════════════════════════
%  front_pages/abstract_nepali.tex
%  Nepali Abstract (शोधसार) — GNN EGFR Thesis
%  Compiled with XeLaTeX
% ════════════════════════════════════════════════════════════════════

\cleardoublepage

\begin{center}
  {\fontsize{18}{22}\selectfont\bfseries{\nepalifont शोधसार}}
\end{center}
\addcontentsline{toc}{chapter}{\texorpdfstring{{\nepalifont शोधसार}}{Shodhasar}}
\thispagestyle{plain}

\vspace{1cm}

{\fontsize{12}{18}\selectfont

% ── Paragraph 1: Background ──────────────────────────────────────────
{\nepalifont फोक्सोको क्यान्सर विश्वव्यापी रूपमा सबैभन्दा बढी
मृत्युदर भएको क्यान्सरहरूमध्ये एक हो, जसमध्ये
नन-स्मल सेल लङ क्यान्सर}
(NSCLC)
{\nepalifont सम्पूर्ण फोक्सोको क्यान्सरको लगभग ८५\% हिस्सा
ओगट्दछ।}
EGFR
{\nepalifont उत्परिवर्तन यस रोगको प्रमुख आणविक चालकको रूपमा
पहिचान गरिएको छ, तथापि अवरोधकहरूप्रति प्राप्त प्रतिरोधक
क्षमता एक प्रमुख चुनौती बनेको छ।}

\vspace{0.3em}

% ── Paragraph 2: Methods & Results ───────────────────────────────────
{\nepalifont यस अध्ययनमा}
EGFR
{\nepalifont अवरोधकहरूको वर्गीकरण र सम्भावित नेतृत्व यौगिकहरू
पहिचान गर्न ग्राफ न्युरल नेटवर्क}
(GNN)
{\nepalifont आधारित गणनात्मक ढाँचा प्रस्तुत गरिएको छ।}
ChEMBL
{\nepalifont डेटाबेसबाट प्रयोगात्मक रूपमा प्रमाणित}
EGFR
{\nepalifont जैव-क्रियाशीलता भएका १०,५२३ यौगिकहरू संकलन गरी,}
RDKit
{\nepalifont प्रयोग गरी स्तरीकृत गरिए र आणविक ग्राफको रूपमा
प्रस्तुत गरियो।}
Bemis-Murcko
{\nepalifont स्क्याफोल्ड विभाजन अन्तर्गत पाँच}
GNN
{\nepalifont वास्तुकला}
(GCN, GAT, GIN, AttentiveFP, GraphSAGE)
{\nepalifont र दुई आधारभूत मोडेल}
(Random Forest, MLP)
{\nepalifont गरी सात मोडेलको तुलना गरियो।}
{\nepalifont फिंगरप्रिन्ट-आधारित आधारभूत मोडेलहरू}
(Random Forest
{\nepalifont AUROC}
0.8827,
MLP
0.8547)
{\nepalifont ले सर्वोच्च प्रदर्शन देखाए, जबकि}
GNN
{\nepalifont हरू}
(0.8087--0.8454,
{\nepalifont जसमध्ये}
GCN
{\nepalifont सबैभन्दा बलियो})
{\nepalifont तुलनायोग्य रहे।}
Optuna
{\nepalifont मार्फत अनुकूलन गरिएको}
AttentiveFP
{\nepalifont ले}
0.8632
{\nepalifont को मान्यकरण}
AUROC
{\nepalifont प्राप्त गर्\kern-0.05em यो, तर यसले आधारभूत मोडेलहरूलाई
उछिन्न सकेन।}

\vspace{0.3em}

% ── Paragraph 3: Nepali plant screening & Conclusion ─────────────────
{\nepalifont अध्ययनका दुई उत्कृष्ट मोडेल (उत्कृष्ट ग्राफ नेटवर्क}
GCN
{\nepalifont र उत्कृष्ट वर्गीकरणकर्ता}
Random Forest)
{\nepalifont लाई}
COCONUT 2.0
{\nepalifont डेटाबेसमार्फत सात परम्परागत नेपाली औषधीय वनस्पति
प्रजातिहरू}
(Swertia chirayita, Rhododendron arboreum, Nardostachys
jatamansi, Terminalia chebula, Berberis aristata, Ocimum sanctum,
Tinospora cordifolia)
{\nepalifont बाट संकलित ५३५ प्राकृतिक यौगिकहरूमध्ये,}
ChEMBL
{\nepalifont तालिम डेटामा पहिले नै रहेका ८ यौगिक (}apigenin
{\nepalifont र}
kaempferol
{\nepalifont सहित) हटाइए, र बाँकी ५२७ वास्तवमै नदेखिएका यौगिकको
स्क्रिनिङ गरियो। यौगिकहरूलाई पूर्वानुमानित सक्रियताका आधारमा
क्रमबद्ध गरी}
Lipinski
{\nepalifont नियम र}
ADMET
{\nepalifont गुणहरू}
(QED, TPSA)
{\nepalifont अनुसार मूल्याङ्कन गरियो। स्क्रिनिङको नतिजा मोडेलमा
निर्भर रह्यो---}
Random Forest
{\nepalifont ले १७९ र}
GCN
{\nepalifont ले २४ सम्भावित-सक्रिय यौगिक पहिचान गरे, तर कुनै पनि
उत्कृष्ट मोडेलले}
p $\geq$ 0.80
{\nepalifont मा}
Lipinski
{\nepalifont-अनुपालित यौगिक दिएन। सबैभन्दा बलियो सम्भावित यौगिकहरू
xanthone}
isogentisin
(Random Forest 0.728)
{\nepalifont र flavone}
luteolin
{\nepalifont (दुवै मोडेलको सहमति) रहे, दुवै मध्यम-विश्वासका।}
{\nepalifont यस अध्ययनले}
GNN
{\nepalifont आधारित गहन अधिगम विधिले}
NSCLC
{\nepalifont उपचारका लागि नेपाली परम्परागत औषधीय वनस्पतिहरूबाट
नयाँ नेतृत्व यौगिकहरू पहिचान गर्न प्रभावकारी भूमिका खेल्न सक्ने
देखाएको छ।}

} % end fontsize

\vspace{0.8cm}

% ── Keywords ─────────────────────────────────────────────────────────
\noindent
{\fontsize{12}{18}\selectfont
{\nepalifont\textbf{मुख्य शब्दहरू:}}
\textit{EGFR,
graph neural networks, AttentiveFP,
virtual screening, Nepali medicinal plants,
natural products, Bemis-Murcko scaffold split,
non-small cell lung cancer, drug discovery}}